\chapter{1978:Daniel Gray Quillen}\label{ch1978danielg.quillen:daniel-gray-quillen-fields-medal-1978-a-comprehensive-report-on-his-work-and-impact}

\begin{center}\rule{0.5\linewidth}{0.5pt}\end{center}

\textbf{Prepared:} July 2026 \textbf{Subject:} Daniel Gray Quillen (b.\
June 27, 1940, Orange, New Jersey, USA; d.~April 30, 2011)
\textbf{Position:} Professor, University of Oxford; MIT
\textbf{Doctoral Advisor:} Raoul Bott (Harvard University)
\textbf{Fields Medal Citation:} \emph{``For his formulation of higher
algebraic K-theory, which has provided a framework for arithmetic and
geometry.''}

\begin{center}\rule{0.5\linewidth}{0.5pt}\end{center}

\begin{prereqs} Algebraic topology (homotopy groups, classifying spaces,
simplicial sets), abstract algebra (groups, rings, modules), category
theory (functors, natural transformations), commutative algebra
(projective modules, localization). \end{prereqs}

\begin{objectives} Construct the plus-construction and Q-construction
for algebraic K-theory; prove the calculation of K-groups of finite
fields; explain the relationship between K-theory and cyclic homology
via the Loday--Quillen--Tsygan theorem; understand Quillen's rational
homotopy theory.
\end{objectives}

\section{Problem History}

Algebraic K-theory originated with Grothendieck's $K_0(X)$ for his
generalized Riemann--Roch theorem. For a ring $A$, $K_0(A)$ is the
Grothendieck group of finitely generated projective modules. Bass
defined $K_1(A) = GL(A)/[GL(A),GL(A)]$ and Milnor defined $K_2(A)$ as
the kernel of the Steinberg group map, but defining $K_n(A)$ for
$n \ge 3$ was the fundamental difficulty.

Quillen's revolutionary insight was to shift from algebra to homotopy
theory. The plus-construction takes a CW complex $X$ and a perfect
normal subgroup $N \unlhd \pi_1(X)$, producing $X^+$ that kills $N$
while inducing homology isomorphisms. Applied to $X = BGL(A)$ and
$N = E(A)$, the space $BGL(A)^+$ has
$K_i(A) = \pi_i(BGL(A)^+)$ for $i \ge 1$.

Quillen computed K-groups of finite fields: for $\mathbb{F}_q$ with
$q = p^d$,
\[
K_{2i-1}(\mathbb{F}_q) \cong \mathbb{Z}/(q^i - 1)\mathbb{Z}, \quad
K_{2i}(\mathbb{F}_q) = 0, \quad i \ge 1,
\]
showing higher K-groups are arithmetically meaningful and computable.

\section{How Was It Solved}

The plus-construction is difficult for schemes. Quillen introduced the
Q-construction for an exact category $\mathcal{C}$: a new category
$Q\mathcal{C}$ where a morphism $M \to N$ is an equivalence class of
diagrams
\[
M \xleftarrow{\text{defl}} Y \xrightarrow{\text{infl}} N,
\]
with $Y$ an admissible epimorphism onto $M$ and admissible monomorphism
into $N$. The K-groups are $K_i(\mathcal{C}) = \pi_{i+1}(BQ\mathcal{C})$.
Quillen proved $+ = Q$: for projective modules,
$\pi_{i+1}(BQ\mathcal{P}(A)) \cong \pi_i(BGL(A)^+)$.

The Q-construction enabled major structural theorems: Devissage
($K_i(\mathcal{A}) \cong K_i(\mathcal{B})$ for filtered abelian
categories) and Localization (a long exact sequence for Serre
subcategories).

Quillen also developed model categories (1967)---categories with weak
equivalences, fibrations, and cofibrations satisfying five axioms---the
standard language for homotopy theory. In rational homotopy theory
(1969), he showed the rational homotopy category of simply connected
spaces is equivalent to differential graded Lie algebras over
$\mathbb{Q}$.

The Loday--Quillen--Tsygan theorem (1984) connects K-theory to cyclic
homology: $H_*(\mathfrak{gl}(A); k) \cong \mathrm{Sym}(HC_{*-1}(A))$,
providing the Chern character map.

\section{Impact of This Discovery}

The Quillen--Lichtenbaum conjecture (largely proved by Voevodsky) relates
K-groups to \'{e}tale cohomology. The Quillen--Suslin theorem settled
Serre's conjecture on projective modules over polynomial rings.

Quillen's discovery that $\pi_*(MU) \cong$ Lazard ring initiated
chromatic homotopy theory. The Quillen metric on determinant line bundles
is key in Arakelov geometry and string theory. In noncommutative geometry,
Quillen's cyclic homology provided tools for Connes's index theory on
noncommutative spaces.

\section{Modern Views and Improvements}

Quillen's model categories were refined into $\infty$-categories (Joyal,
Lurie), now the standard language for higher algebra. Derived algebraic
geometry (To\"{e}n, Vezzosi, Lurie) replaces commutative rings with
simplicial commutative rings using Quillen's cotangent complex.

Motivic homotopy theory (Morel--Voevodsky) applies model categories to
schemes. Topological cyclic homology (TC) approximates K-theory
computationally; the resolution of the redshift conjecture
(Mathew--Naumann--Noel) revealed deep connections with chromatic homotopy.

\begin{center}\rule{0.5\linewidth}{0.5pt}\end{center}

\section*{Summary}
\addcontentsline{toc}{section}{Summary}

Quillen defined higher algebraic K-theory via the plus-construction and
Q-construction, proved the K-groups of finite fields, developed model
categories and rational homotopy theory, and connected K-theory to cyclic
homology, unifying arithmetic, geometry, and topology.

\begin{center}\rule{0.5\linewidth}{0.5pt}\end{center}
